\documentclass[11pt]{article}
% Uses the canonical latex-docs semantic style in-repository. The local
% lightweight fallback is for independent compilation of this distribution.
\IfFileExists{technical-user-manual.sty}{\usepackage{technical-user-manual}}{\usepackage{../support/aaoa-standalone}}
\usepackage{mathtools}
\setDocTitle{Chapter 1: Analysis of Algorithms}
\setDocSubtitle{Guided reading, derivations and original exercises}
\setDocVersion{1.0}
\setDocStatus{Educational reference / unapproved working material}
\setDocOwner{Algorithm Engineering}
\setDocAuthor{Original companion material}
\setDocSummary{Coverage of primary book Chapter 1 with independently authored examples, proof obligations and lab transfer.}
\providecommand{\Expect}{\mathbb{E}}
\providecommand{\Prob}{\mathbb{P}}
\providecommand{\Var}{\operatorname{Var}}
\providecommand{\bigoh}{\mathcal{O}}
\begin{document}
\makedoctocpage
\section{Reading map and prerequisite check}
\noindent\textbf{Primary source sections:} 1.1 Why Analyze an Algorithm?; 1.2 Theory of Algorithms; 1.3 Analysis of Algorithms; 1.4 Average-Case Analysis; 1.5 Example: Analysis of Quicksort; 1.6 Asymptotic Approximations; 1.7 Distributions; 1.8 Randomized Algorithms (printed pp. 3--33).\par
\noindent\textbf{Prerequisites:} Proof notation, logarithms, probability and basic sorting.\par
\section{Learning outcomes}
\begin{enumerate}
\item Distinguish worst-case guarantees, probabilistic expected costs, and machine-specific timing.
\item Create a uniform-random-permutation model and identify where it is unrealistic.
\item Derive and instrument a recurrence for quicksort comparisons.
\end{enumerate}
\section{Guided technical analysis}
\subsection{Two complementary questions}
A complexity bound tells which scales are possible; a cost model tells what a specific algorithm is expected to do. Quicksort has quadratic worst-case comparisons under poor pivot choices, yet its random-permutation average is $\Theta(n\log n)$. Do not describe a single sample as an average-case proof.
\subsection{Count an operation before timing it}
Under partitioning that compares the pivot to every other key, the expected number $Q_n$ of comparisons for a random pivot rank satisfies $Q_0=Q_1=0$ and
\[
Q_n=(n-1)+\frac{2}{n}\sum_{k=0}^{n-1}Q_k.
\]
By eliminating the sum, $nQ_n=(n+1)Q_{n-1}+2(n-1)$ for $n\ge2$. The solution is $Q_n=2(n+1)H_n-4n$. The constants depend on the partition scheme and exact comparison-count definition; do not universalize this identity to every quicksort implementation.
\section{Independent problem set}
\begin{enumerate}
\item Prove the displayed quicksort recurrence from a uniform pivot-rank assumption and enumerate $n=2,3,4$.
\item Construct two pivot rules and one adversarial input family for each; compare comparison counts.
\item Specify how duplicate keys and three-way partitioning change the experiment and model.
\end{enumerate}
\section{Laboratory transfer}
\textbf{Required laboratory:} LAB-PY01: random-permutation quicksort comparisons and small-$n$ oracle.\par
\noindent\textbf{Minimum recorded evidence:} repository commit or source hash; test inputs/seeds; actual operation counts; predicted values; one explicit model limitation.\par
\section{Mastery gate and self-review}
Can reproduce the $n=3$ expected count by enumerating all permutations; can explain why an input distribution is an assumption, not an implementation property.\par
\begin{enumerate}\item Can I define the random object/input exactly? \item Is my identity exact, asymptotic, or empirical? \item Which base cases and counterexamples did I check? \item What would falsify my performance hypothesis?\end{enumerate}
\section{Suggested notebook structure}
Use five named sections: model and invariants; derivation; deterministic small instances; randomized replication; discrepancy and revised hypothesis. Keep a proof separate from measurement results.

\end{document}
